\chapter{Arrow's Impossibility Theorem}
%%%%%%%%%%%%%%%%%%%%%%%%%%%%%%%%%%%%%%%%%%%%%%%%%%%%%%%%%%%%%%%%%%%%%%%%%%%%%%%
\section{Overview}

This chapter proves \textbf{Arrow's Impossibility Theorem} via Kalai's
Fourier-analytic approach.  For 3 alternatives $\{a,b,c\}$ and $n$ voters,
a social welfare function is modelled by an odd Boolean function
$f:\{0,1\}^n\to\bbr$ applied to each pairwise comparison.  The main result is:

\begin{quote}
Any odd, $\pm 1$-valued, unanimous, acyclic social welfare function must be a
\emph{dictator}.
\end{quote}

The proof proceeds in three steps:
\begin{enumerate}
  \item Acyclicity forces the Fourier correlation function $\mathrm{corr}(f) = -1/3$.
  \item Since $\mathrm{corr}(f)\ge -1/3$ always holds (with equality iff all
        Fourier weight is on level 1), this forces $W_1[f] = 1$.
  \item A $\pm 1$-valued degree-1 unanimous function must be a dictator.
\end{enumerate}

\textbf{References:}
\begin{itemize}
  \item Gil Kalai, \emph{A Fourier-Theoretic Perspective on the Condorcet
        Paradox and Arrow's Theorem}, Advances in Applied Mathematics, 2002.
  \item Ryan O'Donnell, \emph{Analysis of Boolean Functions}, Chapter 2.
\end{itemize}

%%%%%%%%%%%%%%%%%%%%%%%%%%%%%%%%%%%%%%%%%%%%%%%%%%%%%%%%%%%%%%%%%%%%%%%%%%%%%%%
\section{Social choice definitions}

\begin{definition}[Unanimity]
\lean{ArrowTheorem.unanimity}
\label{ArrowTheorem.unanimity}
\leanok
\uses{BooleanAnalysis.BooleanFunc}
A social welfare function $f:\{0,1\}^n\to\bbr$ is \emph{unanimous} if
$f(\texttt{false},\dots,\texttt{false}) = 1$: when all voters prefer
alternative $a$ over $b$, so does society.
\end{definition}

\begin{definition}[Dictatorship]
\lean{ArrowTheorem.isDictator}
\label{ArrowTheorem.isDictator}
\leanok
\uses{BooleanAnalysis.dictator}
$f$ is a \emph{dictatorship} if there exists a voter $i$ such that
$f = \mathrm{dict}_i$, i.e.\ society's preference always equals voter $i$'s.
\end{definition}

%%%%%%%%%%%%%%%%%%%%%%%%%%%%%%%%%%%%%%%%%%%%%%%%%%%%%%%%%%%%%%%%%%%%%%%%%%%%%%%
\section{Voter ordering model}

Each voter holds one of the 6 strict transitive orderings of $\{a,b,c\}$.
For each ordering we record three pairwise preferences using the sign
convention $\sigma(\texttt{false}) = 1$ (``pro first alternative'') and
$\sigma(\texttt{true}) = -1$ (``pro second'').

The six orderings and their signs $(s_{ab}, s_{bc}, s_{ca})$ are:
\begin{center}
\begin{tabular}{clccc}
Index & Ordering & $s_{ab}$ & $s_{bc}$ & $s_{ca}$ \\\hline
0 & $a>b>c$ & $1$ & $1$ & $-1$ \\
1 & $a>c>b$ & $1$ & $-1$ & $-1$ \\
2 & $b>a>c$ & $-1$ & $1$ & $-1$ \\
3 & $b>c>a$ & $-1$ & $1$ & $1$ \\
4 & $c>a>b$ & $1$ & $-1$ & $1$ \\
5 & $c>b>a$ & $-1$ & $-1$ & $1$ \\
\end{tabular}
\end{center}

\begin{definition}[Pairwise preference functions]
\lean{ArrowTheorem.abPref}
\label{ArrowTheorem.abPref}
\lean{ArrowTheorem.bcPref}
\label{ArrowTheorem.bcPref}
\lean{ArrowTheorem.caPref}
\label{ArrowTheorem.caPref}
\leanok
\uses{BooleanAnalysis.boolToSign}
$\mathrm{abPref}(k)$, $\mathrm{bcPref}(k)$, $\mathrm{caPref}(k)$ give the
Boolean preference of ordering $k\in\{0,\dots,5\}$ in the $a$-vs-$b$,
$b$-vs-$c$, and $c$-vs-$a$ comparisons, respectively.
\end{definition}

\begin{definition}[Profile]
\lean{ArrowTheorem.Profile}
\label{ArrowTheorem.Profile}
\leanok
A \emph{profile} assigns each of the $n$ voters one of the 6 orderings:
$p : \mathrm{Fin}\,n \to \mathrm{Fin}\,6$.
\end{definition}

\begin{definition}[Vote vectors]
\lean{ArrowTheorem.abVotes}
\label{ArrowTheorem.abVotes}
\lean{ArrowTheorem.bcVotes}
\label{ArrowTheorem.bcVotes}
\lean{ArrowTheorem.caVotes}
\label{ArrowTheorem.caVotes}
\leanok
\uses{ArrowTheorem.Profile, ArrowTheorem.abPref, ArrowTheorem.bcPref, ArrowTheorem.caPref}
Given a profile $p$, the vote vectors $\mathrm{abVotes}(p)$,
$\mathrm{bcVotes}(p)$, $\mathrm{caVotes}(p) \in \{0,1\}^n$ record each
voter's pairwise preference.
\end{definition}

\begin{definition}[Acyclicity]
\lean{ArrowTheorem.acyclic}
\label{ArrowTheorem.acyclic}
\leanok
\uses{ArrowTheorem.abVotes, ArrowTheorem.bcVotes, ArrowTheorem.caVotes}
A social welfare function $f$ is \emph{acyclic} if no profile of transitive
voter orderings produces a Condorcet cycle in society's preferences.
A cycle occurs when
$f(\mathrm{abVotes}(p)) = f(\mathrm{bcVotes}(p)) = f(\mathrm{caVotes}(p)) = 1$
(the cycle $a>b>c>a$) or all equal $-1$ (the reverse cycle).
\end{definition}

%%%%%%%%%%%%%%%%%%%%%%%%%%%%%%%%%%%%%%%%%%%%%%%%%%%%%%%%%%%%%%%%%%%%%%%%%%%%%%%
\section{Key correlation lemmas}

The following lemma encodes the core computation: for a single voter drawn
uniformly from the 6 orderings, the expected product of any two distinct
pairwise preference signs is $-1/3$.

\begin{lemma}[Sum of paired preference signs]
\lean{ArrowTheorem.sum_abPref_bcPref}
\label{ArrowTheorem.sum_abPref_bcPref}
\statementsource{boolean-ch02-social-choice-arrow}{pdf-352ab7ff3113-p017-b002}
\lean{ArrowTheorem.sum_bcPref_caPref}
\label{ArrowTheorem.sum_bcPref_caPref}
\proofsource{boolean-ch02-social-choice-arrow}{
  pdf-352ab7ff3113-p017-b002,
  pdf-352ab7ff3113-p016-b001
}
\lean{ArrowTheorem.sum_abPref_caPref}
\label{ArrowTheorem.sum_abPref_caPref}
\leanok
\uses{ArrowTheorem.abPref, ArrowTheorem.bcPref, ArrowTheorem.caPref, BooleanAnalysis.boolToSign}
\difficulty{1}
Label the six strict orderings of three alternatives $a,b,c$ by $k\in\{0,\dots,5\}$. For
each ordering let $x_k\in\{-1,+1\}$ record the $a$-versus-$b$ comparison, taking the
value $+1$ when $a$ is preferred to $b$ and $-1$ when $b$ is preferred to $a$, and let
$y_k\in\{-1,+1\}$ record the $b$-versus-$c$ comparison in the same way, with $+1$ when
$b$ is preferred to $c$ and $-1$ otherwise. Then
\[
\sum_{k=0}^{5} x_k\,y_k \;=\; -2.
\]

For each of the six strict orderings $k \in \{0,\dots,5\}$ of three alternatives
$a,b,c$, let $\sigma$ be the sign encoding, sending $\mathrm{false}$ to $+1$ and
$\mathrm{true}$ to $-1$. Write $y_k = \sigma(p_k)$, where $p_k$ is the Boolean
preference of ordering $k$ in the $b$-vs-$c$ comparison, and $z_k = \sigma(q_k)$, where
$q_k$ is the preference in the $c$-vs-$a$ comparison. Then
\[
\sum_{k=0}^{5} y_k\, z_k \;=\; -2.
\]

For each of the six strict orderings $k \in \{0,\dots,5\}$ of three alternatives
$a,b,c$, write $p_{ab}(k)$ and $p_{ca}(k)$ for the Boolean preferences of ordering $k$
in the $a$-vs-$b$ and $c$-vs-$a$ comparisons, and let $\sigma$ be the sign encoding
sending $\mathrm{false}$ to $1$ and $\mathrm{true}$ to $-1$. Then the sum over all six
orderings of the products of the encoded signs is
\[
\sum_{k=0}^{5} \sigma\big(p_{ab}(k)\big)\,\sigma\big(p_{ca}(k)\big)
\;=\; -2.
\]
\end{lemma}

%%%%%%%%%%%%%%%%%%%%%%%%%%%%%%%%%%%%%%%%%%%%%%%%%%%%%%%%%%%%%%%%%%%%%%%%%%%%%%%
\section{Fourier formula for the cycle correlation}

\begin{definition}[Correlation function]
\lean{ArrowTheorem.corrFunc}
\label{ArrowTheorem.corrFunc}
\leanok
\uses{BooleanAnalysis.fourierCoeff}
The \emph{Fourier correlation function} of $f$ is
\[
  \mathrm{corr}(f) \;=\; \sum_{S\subseteq[n]} \hat f(S)^2\,(-1/3)^{|S|}.
\]
For odd $f$, only odd-level terms contribute.
\end{definition}

\begin{lemma}[Expected pairwise vote product equals the Fourier correlation function]
\lean{ArrowTheorem.expected_product_eq_corrFunc}
\label{ArrowTheorem.expected_product_eq_corrFunc}
\leanok
\proofsource{boolean-ch02-social-choice-arrow}{pdf-352ab7ff3113-p017-b002}
\uses{ArrowTheorem.corrFunc, ArrowTheorem.abVotes, ArrowTheorem.bcVotes, BooleanAnalysis.fourierCoeff, BooleanAnalysis.chiS, ArrowTheorem.sum_abPref_bcPref}
\difficulty{1}
Let $f:\{0,1\}^n\to\bbr$ be a Boolean function. Suppose each of the $n$ voters chooses,
independently and uniformly, one of the six strict orderings of three candidates
$a,b,c$; for such a profile $p$ write $x_p\in\{0,1\}^n$ for the vector whose $i$-th
entry records voter $i$'s preference in the $a$-versus-$b$ comparison, and
$y_p\in\{0,1\}^n$ for the vector recording each voter's preference in the $b$-versus-$c$
comparison. Then, averaging the product of the values of $f$ on these two vectors over
all $6^n$ profiles,
\[
  \frac{1}{6^{\,n}}\sum_{p}\, f(x_p)\,f(y_p)
  \;=\; \sum_{S\subseteq[n]} \hat f(S)^2\,(-1/3)^{|S|},
\]
where $\hat f(S)$ denotes the Fourier--Walsh coefficient of $f$ at frequency $S$.
\end{lemma}

\begin{lemma}[Lower bound for the Fourier correlation function]
\lean{ArrowTheorem.corrFunc_ge_neg_third}
\label{ArrowTheorem.corrFunc_ge_neg_third}
\leanok
\proofsource{boolean-ch02-social-choice-arrow}{pdf-352ab7ff3113-p017-b003}
\uses{ArrowTheorem.corrFunc, BooleanAnalysis.isOddFunc, BooleanAnalysis.isPmOne, BooleanAnalysis.fourierCoeff_odd_even, BooleanAnalysis.parseval_pm_one}
\difficulty{5}
Let $f:\{0,1\}^n\to\bbr$ be a Boolean function that is odd, so that $f(\lnot x)=-f(x)$
for all $x$, and $\pm 1$-valued, so that $f(x)\in\{-1,1\}$ for all $x$. Then its Fourier
correlation function satisfies
\[
\mathrm{corr}(f) \;=\; \sum_{S\subseteq[n]} \hat f(S)^2\,(-1/3)^{|S|} \;\ge\; -\tfrac13,
\]
where $\hat f(S)$ denotes the Fourier--Walsh coefficient of $f$ at frequency $S$.
\end{lemma}

\begin{lemma}[Extremal correlation concentrates weight on level one]
\lean{ArrowTheorem.corrFunc_eq_neg_third_of_weight_one}
\label{ArrowTheorem.corrFunc_eq_neg_third_of_weight_one}
\leanok
\proofsource{boolean-ch02-social-choice-arrow}{pdf-352ab7ff3113-p017-b003}
\uses{ArrowTheorem.corrFunc, ArrowTheorem.corrFunc_ge_neg_third, BooleanAnalysis.isOddFunc, BooleanAnalysis.isPmOne, BooleanAnalysis.fourierCoeff_odd_even, BooleanAnalysis.parseval_pm_one}
\difficulty{6}
Let $f:\{0,1\}^n\to\bbr$ be a $\pm 1$-valued Boolean function that is odd, meaning
$f(\lnot x)=-f(x)$ for every $x$, and suppose its Fourier correlation function attains
the value
\[
  \mathrm{corr}(f)\;=\;\sum_{S\subseteq[n]}\hat f(S)^2\,(-1/3)^{|S|}\;=\;-\tfrac13.
\]
Then $\hat f(S)=0$ for every subset $S\subseteq[n]$ with $|S|\neq 1$; that is, all
Fourier--Walsh weight of $f$ is carried by the singleton frequencies.
\end{lemma}

%%%%%%%%%%%%%%%%%%%%%%%%%%%%%%%%%%%%%%%%%%%%%%%%%%%%%%%%%%%%%%%%%%%%%%%%%%%%%%%
\section{Acyclicity implies degree-1 structure}

\begin{lemma}[Acyclicity forces correlation $-1/3$]
\lean{ArrowTheorem.acyclic_implies_corrFunc}
\label{ArrowTheorem.acyclic_implies_corrFunc}
\leanok
\proofsource{boolean-ch02-social-choice-arrow}{
  pdf-352ab7ff3113-p017-b002,
  pdf-352ab7ff3113-p017-b003
}
\uses{ArrowTheorem.acyclic, ArrowTheorem.corrFunc, ArrowTheorem.expected_product_eq_corrFunc, BooleanAnalysis.isPmOne}
\difficulty{5}
Let $f:\{0,1\}^n\to\bbr$ be a Boolean function that is odd, $\pm 1$-valued, and acyclic.
Then its Fourier correlation function satisfies
\[
  \mathrm{corr}(f)\;=\;\sum_{S\subseteq[n]}\hat f(S)^2\,(-1/3)^{|S|}\;=\;-\tfrac{1}{3}.
\]
\end{lemma}

%%%%%%%%%%%%%%%%%%%%%%%%%%%%%%%%%%%%%%%%%%%%%%%%%%%%%%%%%%%%%%%%%%%%%%%%%%%%%%%
\section{Degree-1 implies dictatorship}

\begin{lemma}[Degree-one unanimous voting rules are dictatorships]
\lean{ArrowTheorem.degree_one_implies_dictator}
\label{ArrowTheorem.degree_one_implies_dictator}
\leanok
\statementsource{boolean-ch02-social-choice-arrow}{pdf-352ab7ff3113-p017-b003}
\uses{ArrowTheorem.isDictator, ArrowTheorem.unanimity, BooleanAnalysis.isPmOne, BooleanAnalysis.fourierCoeff, BooleanAnalysis.parseval_pm_one, BooleanAnalysis.walsh_expansion, BooleanAnalysis.dictator_eq_chi, BooleanAnalysis.chiS_singleton}
\difficulty{6}
Let $f:\{0,1\}^n\to\bbr$ be a Boolean function that is odd, is $\pm 1$-valued, and is
unanimous. Suppose further that $f$ has degree one, in the sense that its Fourier--Walsh
coefficient $\hat f(S)$ vanishes for every set $S\subseteq[n]$ with $\abs{S}\neq 1$.
Then $f$ is a dictatorship: there is a coordinate $i$ with $f=\mathrm{dict}_i$.
\end{lemma}

%%%%%%%%%%%%%%%%%%%%%%%%%%%%%%%%%%%%%%%%%%%%%%%%%%%%%%%%%%%%%%%%%%%%%%%%%%%%%%%
\section{Arrow's Impossibility Theorem}

\begin{theorem}[Arrow's impossibility theorem]
\lean{ArrowTheorem.arrow_theorem}
\label{ArrowTheorem.arrow_theorem}
\leanok
\proofsource{boolean-ch02-social-choice-arrow}{
  pdf-352ab7ff3113-p016-b004,
  pdf-352ab7ff3113-p017-b003
}
\uses{ArrowTheorem.acyclic_implies_corrFunc, ArrowTheorem.corrFunc_eq_neg_third_of_weight_one, ArrowTheorem.degree_one_implies_dictator}
\difficulty{2}
Let $n$ be a natural number and let $f:\{0,1\}^n\to\bbr$ be a social welfare function
that aggregates the pairwise opinions of $n$ voters. Suppose $f$ is odd, meaning $f(\bar
x) = -f(x)$ for every $x$, where $\bar x$ is the bitwise complement of $x$; is $\pm
1$-valued, meaning $f(x)\in\{-1,1\}$ for every $x$; is unanimous, meaning $f$ takes the
value $1$ at the all-zeros input; and is acyclic, meaning that no profile of voter
orderings produces a Condorcet cycle. Here a profile assigns to each of the $n$ voters
one of the six linear orderings of three alternatives $a,b,c$, and from it one forms
three vote vectors in $\{0,1\}^n$ whose $i$-th entries record voter $i$'s preference in
the $a$-vs-$b$, $b$-vs-$c$, and $c$-vs-$a$ comparisons; acyclicity requires that $f$
never sends all three of these vectors to $1$, and never sends all three to $-1$. Then
$f$ is a dictatorship: there is a voter $i$ such that $f(x) = (-1)^{x_i}$ for all $x$.
\end{theorem}

%%%%%%%%%%%%%%%%%%%%%%%%%%%%%%%%%%%%%%%%%%%%%%%%%%%%%%%%%%%%%%%%%%%%%%%%%%%%%%%
%%% added by the blueprint pipeline
\section{Additional declarations}

\begin{lemma}[Balance of the a-vs-b preference signs]
\lean{ArrowTheorem.sum_abPref_sign}
\label{ArrowTheorem.sum_abPref_sign}
\leanok
\uses{ArrowTheorem.abPref, BooleanAnalysis.boolToSign}
\difficulty{1}
There are six strict orderings $k \in \{0,1,\dots,5\}$ of the three alternatives
$a,b,c$. For each such ordering, let $p_{ab}(k)$ be its $a$-vs-$b$ preference, the
Boolean that records whether $b$ is ranked above $a$, and apply the sign encoding
$\sigma$, so that $\sigma(p_{ab}(k))$ equals $+1$ when ordering $k$ ranks $a$ above $b$
and $-1$ when it ranks $b$ above $a$. Then these six signed preferences sum to zero:
\[
  \sum_{k=0}^{5} \sigma\bigl(p_{ab}(k)\bigr) = 0.
\]
\end{lemma}

\begin{lemma}[Balanced b-versus-c preferences over the six orderings]
\lean{ArrowTheorem.sum_bcPref_sign}
\label{ArrowTheorem.sum_bcPref_sign}
\leanok
\uses{ArrowTheorem.bcPref, BooleanAnalysis.boolToSign}
\difficulty{1}
Consider the six strict orderings of three alternatives $a$, $b$, and $c$, indexed by $k
\in \{0,\dots,5\}$, and let $p(k)$ denote the $b$-versus-$c$ pairwise preference of
ordering $k$. Under the sign encoding $\sigma$, these preferences sum to zero across all
six orderings:
\[
  \sum_{k=0}^{5} \sigma\bigl(p(k)\bigr) = 0.
\]
\end{lemma}

\begin{lemma}[Vanishing sum of signed $c$-vs-$a$ preferences]
\lean{ArrowTheorem.sum_caPref_sign}
\label{ArrowTheorem.sum_caPref_sign}
\leanok
\uses{ArrowTheorem.caPref, BooleanAnalysis.boolToSign}
\difficulty{1}
For each of the six strict orderings $k\in\{0,1,\dots,5\}$ of three alternatives, let
$b_k\in\mathrm{Bool}$ record the $c$-vs-$a$ preference of ordering $k$, and encode it as
a sign by $\sigma(b_k)=1$ when $b_k$ is false and $\sigma(b_k)=-1$ when $b_k$ is true.
Then these six signed preferences sum to zero:
\[
  \sum_{k=0}^{5} \sigma(b_k) = 0.
\]
\end{lemma}

\begin{lemma}[Product over a finset as an indicator product]
\lean{ArrowTheorem.prod_finset_eq_prod_univ_ite}
\label{ArrowTheorem.prod_finset_eq_prod_univ_ite}
\leanok
\difficulty{2}
Let $A \subseteq \mathrm{Fin}\,n$ be a finite set of indices and let $g :
\mathrm{Fin}\,n \to \bbr$. Then
\[
  \prod_{j \in A} g(j) \;=\; \prod_{j : \mathrm{Fin}\,n}
    \begin{cases} g(j) & \text{if } j \in A, \\ 1 & \text{if } j \notin A, \end{cases}
\]
so that a product taken only over the members of $A$ may equivalently be written as a
product over all of $\mathrm{Fin}\,n$ in which each factor outside $A$ is replaced by
$1$.
\end{lemma}

\begin{lemma}[General profile kernel for paired preference assignments]
\lean{ArrowTheorem.profile_kernel_gen}
\label{ArrowTheorem.profile_kernel_gen}
\leanok
\proofsource{boolean-ch02-social-choice-arrow}{pdf-352ab7ff3113-p017-b002}
\uses{ArrowTheorem.Profile, BooleanAnalysis.boolToSign, BooleanAnalysis.chiS, ArrowTheorem.prod_finset_eq_prod_univ_ite}
\difficulty{6}
Let $x, y : \mathrm{Fin}\,6 \to \mathrm{Bool}$ be two preference assignments whose sign
encodings have balanced marginals, $\sum_{k=1}^{6}\sigma(x_k) =
\sum_{k=1}^{6}\sigma(y_k) = 0$, and whose signed correlation is
$\sum_{k=1}^{6}\sigma(x_k)\,\sigma(y_k) = -2$. For a profile $p : \mathrm{Fin}\,n \to
\mathrm{Fin}\,6$, write $x\circ p$ and $y\circ p$ for the points of the Boolean
hypercube whose $i$-th coordinates are $x_{p(i)}$ and $y_{p(i)}$. Then for all subsets
$S, T \subseteq [n]$, the average of the product of Walsh characters over all $6^n$
profiles satisfies
\[
  \Bigl(\tfrac16\Bigr)^{n}\sum_{p:\mathrm{Fin}\,n \to \mathrm{Fin}\,6}
    \chi_S(x\circ p)\,\chi_T(y\circ p)
  \;=\;
  \begin{cases} \bigl(-\tfrac13\bigr)^{\abs{S}} & S = T,\\[2pt] 0 & S \ne T. \end{cases}
\]
\end{lemma}

\begin{lemma}[Orthogonality kernel for the $ab$–$bc$ vote pair]
\lean{ArrowTheorem.profile_inner_product_kernel}
\label{ArrowTheorem.profile_inner_product_kernel}
\leanok
\proofsource{boolean-ch02-social-choice-arrow}{
  pdf-352ab7ff3113-p017-b002,
  pdf-352ab7ff3113-p016-b001
}
\uses{ArrowTheorem.Profile, ArrowTheorem.abVotes, ArrowTheorem.bcVotes, ArrowTheorem.sum_abPref_bcPref, BooleanAnalysis.chiS, ArrowTheorem.profile_kernel_gen, ArrowTheorem.sum_abPref_sign, ArrowTheorem.sum_bcPref_sign}
\difficulty{1}
Fix $n$, and for each profile $p \in \{0,\dots,5\}^n$ let $v_{ab}(p), v_{bc}(p) \in
\{0,1\}^n$ be the vote vectors recording, voter by voter, each voter's pairwise
preference in the $a$-vs-$b$ and $b$-vs-$c$ comparisons, and for $S \subseteq [n]$ let
$\chi_S$ be the associated Walsh–Fourier character. Then for all subsets $S, T \subseteq
[n]$, the uniform average of the product $\chi_S(v_{ab}(p))\,\chi_T(v_{bc}(p))$ over the
$6^n$ profiles satisfies
\[
  \Bigl(\tfrac{1}{6}\Bigr)^{\!n}\sum_{p}
    \chi_S\bigl(v_{ab}(p)\bigr)\,\chi_T\bigl(v_{bc}(p)\bigr)
  \;=\;
\begin{cases} \left(-\tfrac{1}{3}\right)^{\abs{S}} & S = T,\\[2pt] 0 & S \ne T.
\end{cases}
\]
\end{lemma}

\begin{lemma}[The $bc$–$ca$ profile kernel]
\lean{ArrowTheorem.profile_kernel_bcca}
\label{ArrowTheorem.profile_kernel_bcca}
\leanok
\proofsource{boolean-ch02-social-choice-arrow}{
  pdf-352ab7ff3113-p016-b001,
  pdf-352ab7ff3113-p017-b002
}
\uses{ArrowTheorem.Profile, ArrowTheorem.bcVotes, ArrowTheorem.caVotes, ArrowTheorem.sum_bcPref_caPref, BooleanAnalysis.chiS, ArrowTheorem.profile_kernel_gen, ArrowTheorem.sum_bcPref_sign, ArrowTheorem.sum_caPref_sign}
\difficulty{1}
Let $S, T \subseteq [n]$, and for a profile $p \colon \mathrm{Fin}\,n \to
\mathrm{Fin}\,6$ write $x(p)$ for its $bc$-vote vector and $y(p)$ for its $ca$-vote
vector, each lying in $\{0,1\}^n$. Averaging the product of Walsh characters over all
$6^n$ profiles gives
\[
  \left(\tfrac{1}{6}\right)^{n} \sum_{p} \chi_S\bigl(x(p)\bigr)\,\chi_T\bigl(y(p)\bigr)
  \;=\;
\begin{cases} \left(-\tfrac{1}{3}\right)^{\abs{S}} & \text{if } S = T,\\[2pt] 0 &
\text{if } S \ne T. \end{cases}
\]
\end{lemma}

\begin{lemma}[Kernel of the ab–ca vote correlation]
\lean{ArrowTheorem.profile_kernel_abca}
\label{ArrowTheorem.profile_kernel_abca}
\leanok
\proofsource{boolean-ch02-social-choice-arrow}{pdf-352ab7ff3113-p017-b002}
\uses{ArrowTheorem.Profile, ArrowTheorem.abVotes, ArrowTheorem.caVotes, ArrowTheorem.sum_abPref_caPref, BooleanAnalysis.chiS, ArrowTheorem.profile_kernel_gen, ArrowTheorem.sum_abPref_sign, ArrowTheorem.sum_caPref_sign}
\difficulty{1}
Fix a natural number $n$ and consider profiles $p$ of $n$ voters, each of whom reports
one of the six strict orderings of the three alternatives $a, b, c$; there are $6^n$
such profiles. For a profile $p$, let $x^{ab}(p) \in \{0,1\}^n$ be the vector recording
each voter's preference in the $a$-vs-$b$ comparison, and $x^{ca}(p) \in \{0,1\}^n$ the
vector recording each voter's preference in the $c$-vs-$a$ comparison. Then for any two
subsets $S, T \subseteq \{1,\dots,n\}$, the averaged product of the corresponding
Walsh–Fourier characters satisfies
\[
\left(\tfrac{1}{6}\right)^{\!n} \sum_{p}
\chi_S\bigl(x^{ab}(p)\bigr)\,\chi_T\bigl(x^{ca}(p)\bigr)
  \;=\;
\begin{cases} \left(-\tfrac{1}{3}\right)^{|S|} & \text{if } S = T,\\[2pt] 0 & \text{if }
S \neq T, \end{cases}
\]
where the sum ranges over all $6^n$ profiles.
\end{lemma}

\begin{lemma}[Expected product from the character kernel]
\lean{ArrowTheorem.expected_product_helper}
\label{ArrowTheorem.expected_product_helper}
\leanok
\proofsource{boolean-ch02-social-choice-arrow}{pdf-352ab7ff3113-p017-b002}
\uses{ArrowTheorem.Profile, ArrowTheorem.corrFunc, BooleanAnalysis.BoolCube, BooleanAnalysis.BooleanFunc, BooleanAnalysis.chiS, BooleanAnalysis.fourierCoeff, BooleanAnalysis.walsh_expansion}
\difficulty{5}
Let $f:\{0,1\}^n\to\bbr$ be a Boolean function, and let
$\mathrm{votes}_1,\mathrm{votes}_2$ each assign to every profile
$p\in(\mathrm{Fin}\,6)^{\,n}$ a point of the Boolean hypercube $\{0,1\}^n$. Suppose the
pair satisfies the kernel identity
\[
\Bigl(\tfrac16\Bigr)^{\!n}\sum_{p}
\chi_S(\mathrm{votes}_1(p))\,\chi_T(\mathrm{votes}_2(p))
  \;=\; \begin{cases}(-1/3)^{|S|} & S = T,\\ 0 & S \ne T,\end{cases}
\]
for all $S,T\subseteq[n]$, where the sum runs over all $6^n$ profiles. Then
\[
  \Bigl(\tfrac16\Bigr)^{\!n}\sum_{p} f(\mathrm{votes}_1(p))\,f(\mathrm{votes}_2(p))
  \;=\; \mathrm{corr}(f)
  \;=\; \sum_{S\subseteq[n]} \hat f(S)^2\,(-1/3)^{|S|}.
\]
\end{lemma}

\begin{lemma}[Expected product of the $bc$ and $ca$ vote vectors]
\lean{ArrowTheorem.expected_product_bcca}
\label{ArrowTheorem.expected_product_bcca}
\leanok
\proofsource{boolean-ch02-social-choice-arrow}{
  pdf-352ab7ff3113-p017-b002,
  pdf-352ab7ff3113-p016-b001
}
\uses{ArrowTheorem.Profile, ArrowTheorem.bcVotes, ArrowTheorem.caVotes, ArrowTheorem.corrFunc, BooleanAnalysis.BooleanFunc, ArrowTheorem.expected_product_helper, ArrowTheorem.profile_kernel_bcca}
\difficulty{1}
Let $f:\{0,1\}^n\to\bbr$ be a Boolean function, and for each profile
$p:\mathrm{Fin}\,n\to\{0,\dots,5\}$ write $x_{bc}(p),\,x_{ca}(p)\in\{0,1\}^n$ for the
$b$-vs-$c$ and $c$-vs-$a$ vote vectors it induces. Then the average of the product
$f\bigl(x_{bc}(p)\bigr)\,f\bigl(x_{ca}(p)\bigr)$ over all $6^n$ profiles equals the
Fourier correlation function of $f$:
\[
  \Bigl(\tfrac16\Bigr)^{n}\sum_{p} f\bigl(x_{bc}(p)\bigr)\,f\bigl(x_{ca}(p)\bigr)
  \;=\; \mathrm{corr}(f)
  \;=\; \sum_{S\subseteq[n]} \hat f(S)^2\,(-1/3)^{|S|}.
\]
\end{lemma}

\begin{lemma}[Expected product of the $ab$ and $ca$ votes]
\lean{ArrowTheorem.expected_product_abca}
\label{ArrowTheorem.expected_product_abca}
\leanok
\proofsource{boolean-ch02-social-choice-arrow}{
  pdf-352ab7ff3113-p017-b002,
  pdf-352ab7ff3113-p016-b001
}
\uses{ArrowTheorem.Profile, ArrowTheorem.abVotes, ArrowTheorem.caVotes, ArrowTheorem.corrFunc, BooleanAnalysis.BooleanFunc, ArrowTheorem.expected_product_helper, ArrowTheorem.profile_kernel_abca}
\difficulty{1}
Let $f:\{0,1\}^n\to\bbr$ be a Boolean function. For a profile $p$ that assigns each of
the $n$ voters one of the six strict orderings of the three alternatives $a,b,c$, let
$x^{ab}(p)\in\{0,1\}^n$ record each voter's preference in the $a$-vs-$b$ comparison and
let $x^{ca}(p)\in\{0,1\}^n$ record each voter's preference in the $c$-vs-$a$ comparison.
Averaging uniformly over all $6^n$ profiles, the expected product of $f$ at these two
vote vectors equals the Fourier correlation function of $f$:
\[
  \Bigl(\tfrac{1}{6}\Bigr)^{n}\sum_{p} f\bigl(x^{ab}(p)\bigr)\,f\bigl(x^{ca}(p)\bigr)
  \;=\; \sum_{S\subseteq[n]} \hat f(S)^{2}\,(-1/3)^{|S|}.
\]
\end{lemma}
